% Abstract: what we did, why it matters, how, the evidence, and the most surprising result.
\begin{abstract}
We test whether putting Coherent Extrapolated Volition (CEV), Yudkowsky's proposal that an AI should do what
humanity would want ``if we knew more, thought faster, were more the people we wished we were'', into a language
model's system prompt makes the model behave in a more aligned way than a classic assistant system prompt.
Alignment theory has proposed such targets for two decades. System prompts are the cheapest way to hand one to a deployed
model, but whether doing so helps, and on which dimensions, has not been measured. We ran four models (GPT, Claude,
Qwen and Llama) under five system prompts: none, classic, CEV, classic+CEV, and a length-matched neutral control. Each
condition was scored on six public evaluations chosen to mirror what CEV is for: human wellbeing, human agency,
manipulation, corrigibility and power-seeking, instrumental convergence, and agentic blackmail. In total there are
\RNSamples{} item-paired samples, graded by an external judge. Relative to the classic prompt, the CEV prompt improved
no evaluation and lowered HumaneBench (\RHbCevClassicDiff, Holm $p=\RHbCevClassicHolm$). Adding the CEV text to the
classic prompt, compared with an equally long neutral text, improved corrigibility and power-seeking answers
(\RAirClassiccevLenmatchDiff, $p=\RAirClassiccevLenmatchHolmAll$ after correction over all primary tests) and lowered
deference to users on HumanAgencyBench (\RHabClassiccevLenmatchDiff, not significant after correction). The losses
concentrate in the \RMentionShareMin--\RMentionShareMax{} of answers that explicitly reason ``as CEV''. These score
\RHabMentionAligned{} on HumanAgencyBench against \RHabNoMentionAligned{} for answers that do not. A judge-blinding
test shows that the judge does not penalise the mention itself. Models that invoke humanity's extrapolated volition
tend to make the user's decisions for them.
\end{abstract}
